\section{Null symmetry and cut-running}\label{sec:null-cut-branch-hinge}

\begin{definition}[Null symmetry]
The AOD form of the AFC Null posture is
\[
\Null\equiv \infD,
\]
where \(\infD\) names unresolved possible distinction. Null names symmetric closure under which branch selection is presented.
\end{definition}

\paragraph{Null chain.}
The Null chain preserves the same syntax as the local cut-running form:
\[
\Null\equiv\infty
\]
\[
\Null\to\Null\to\Null\equiv\Null.
\]

\begin{definition}[Cut operator]
The AOD cut operator sends a retained occurrence into simultaneous identity and cut-running:
\begin{equation}
x\mapsto x\succ x=x.
\label{eq:cut-running-form}
\end{equation}
The symbol \(\succ\) names asymmetric cut-running. The expression records that the same \(x\) is held in identity while also presenting a running cut relation.
\end{definition}

\paragraph{Stokes identity as temporal dynamics.}
The cut operator presents the Stokes identity as temporal dynamics. A null bounds its unspooled infinity of nulls and returns to null. This completed return is the monon cycle:
\begin{equation}
\mathrm{monon}(x)=\operatorname{cycle}(x\succ x=x).
\label{eq:monon-cycle}
\end{equation}
The monon is the first completed beat of the cut operator. It is the primitive completion marker of AOD return dynamics.

\paragraph{Simultaneous presentation.}
The null potential is symmetric in totality. The cut operator unspools that symmetry into temporal asymmetry. Motion is the rebalancing of the unspooled asymmetry toward symmetric closure. The rebalancing has no terminal global end because the null remains simultaneously symmetric while cut-running continues. Local completion is the monon beat: null-to-unspooled-boundary-to-null return.

\paragraph{Inheritance.}
Each occurrence of \(x\) inherits the same primitive presentation:
\[
x\mapsto (x\succ x=x).
\]
Thus
\[
x
\Rightarrow (x\succ x=x)
\Rightarrow ((x\succ x=x)\succ(x\succ x=x)=(x\succ x=x))
\Rightarrow \cdots.
\]
The cut begets a cut because every presented occurrence can present the same simultaneous identity and running relation.

\begin{definition}[Branch / hinge / branch]
A tracked cut-running presentation distinguishes branch, hinge, and branch. The branches are the two retained occurrences exposed by the cut-running relation. The hinge is the middle occurrence through which return, dwell, pinning, reseating, or closure can occur.
\end{definition}

\begin{definition}[Hinge triad]
The branch/hinge/branch support is the hinge triad
\[
\mathsf H_3=(x^-,\mu,x^+).
\]
Here \(x^-\) is the past branch, \(\mu\) is the hinge / dwell / pivot, and \(x^+\) is the future branch. Its local ternary support is
\[
\{-1,0,+1\}.
\]
\end{definition}

\paragraph{Past, present, and future.}
Use the direct temporal terms
\[
x^- = \text{past},
\qquad
x^0 = \text{present},
\qquad
x^+ = \text{future}.
\]
The temporal dynamic is
\[
x^-\to x^0\to x^+.
\]
The field holds past, present, and future in symmetric totality. The cut crowns the present:
\[
\operatorname{Crown}(x)=x^0.
\]
Temporal dynamics is the asymmetric unspooling of this symmetric totality.

\paragraph{Retained ternary support.}
When branch/hinge/branch is retained locally, its coarse support is
\[
\{-1,0,+1\}.
\]
The value \(0\) is hinge, dwell, or pivot. The values \(-1\) and \(+1\) are branch orientations. For \(k\) unconstrained retained positions, the retained base-3 capacity is
\[
C_3(k)=\log_3(3^k)=k.
\]
This capacity is carried by duon current and by entropic duon current when unresolved return burden is retained.
